\begin{table}[htbp]
\caption{Interdisciplinary connections as a research design. The right-hand column states research directions proposed in this chapter, rather than completed experiments.}
\label{tab:ch11-bridges}
\centering\small\renewcommand{\arraystretch}{1.17}
\begin{tabularx}{\linewidth}{@{}>{\raggedright\arraybackslash}p{.24\linewidth}>{\raggedright\arraybackslash}X>{\raggedright\arraybackslash}X@{}}
\toprule
\textbf{Connection} & \textbf{Contribution to open source evaluation} & \textbf{Reciprocal research opportunity}\\\midrule
Measurement theory & Distinguish an index, an attribute, its interpretation, and the consequences of use. & Study interpretation and measurement change when a public score becomes an incentive.\\[6pt]
Network science & Match relations and centrality to the process being represented; distinguish position from causal influence. & Compare graph constructions against documented episodes of open collaboration.\\[6pt]
Software engineering and visual analytics & Connect events with work practices; validate domain tasks, representations, interfaces, and algorithms separately. & Assess whether tools help users reconstruct and correctly interpret evaluation evidence.\\[6pt]
Education and learning analytics & Relate participation to learning objectives, work samples, feedback, and opportunities to learn. & Follow transitions from course participation to sustained community contribution.\\[6pt]
Institutional research and science policy & Explain resource conditions, collective action, and the behavioural effects of public evaluation. & Examine how recognition and decision-support arrangements alter collaboration and maintenance.\\
\bottomrule
\end{tabularx}
\tablenote{The connections draw on the primary literature discussed in Section~11.5. They preserve differences among theoretical positions and avoid treating a borrowed concept as independent evidence for a claim about open source.}
\end{table}
